% chapters/08/section-01.tex
\section{Dari yang Digerakkan Menjadi yang Bergerak Sendiri}
\label{sec:bergerak-sendiri}

Sampai bab sebelumnya, kecepatan setiap benda ditentukan dari luar dirinya.
Bagian ini menjelaskan apa yang berubah ketika benda mulai menentukan
kecepatannya sendiri.

\subsection{Siapa yang Menentukan Kecepatan}

Perhatikan asal kecepatan pada tiga bab terakhir.

\begin{itemize}
  \item \textbf{Pada Bab~V}, kecepatan pemain berasal dari tombol yang ditekan.
        Bendanya pasif: ia bergerak hanya karena diperintah.
  \item \textbf{Pada Bab~VI}, kecepatan tegak berubah karena percepatan
        gravitasi. Perubahannya tetap dan sama untuk semua benda; tidak ada
        keputusan di dalamnya.
  \item \textbf{Pada Bab~VII}, kecepatan berubah karena tabrakan. Perubahan itu
        juga bukan keputusan, melainkan akibat dari keadaan.
\end{itemize}

Pada bab ini, kecepatan berasal dari \textbf{keputusan benda itu sendiri}
berdasarkan keadaan di sekitarnya. Bendanya menjadi agen, yaitu sesuatu yang
mengamati, memutuskan, dan bertindak. Perbedaan itu tampak pada kode: fungsi
\code{update} benda tidak lagi menerima arah dari luar, melainkan menghitung
arahnya sendiri dari posisi pemain.

\begin{catatabox}
Pandangan permainan berubah pada bab ini. Bab~VI dan Bab~VII
memakai pandangan samping, sehingga gravitasi menjadi unsur utama. Bab ini
memakai \textbf{pandangan atas} (dilihat dari langit-langit), sehingga gravitasi
tidak diperlukan dan yang menjadi unsur utama adalah arah pada bidang datar.
Pilihan pandangan itu merupakan keputusan perancangan, bukan keharusan teknis.
\end{catatabox}

\subsection{Susunan Keputusan sebuah Agen}

Bagian~2.5 telah menjelaskan bahwa pengendalian benda memerlukan dua lapisan,
yaitu lapisan keputusan dan lapisan navigasi. Bab ini menyusun lapisan
keputusan yang paling sederhana, yaitu \textbf{perilaku reaktif}: keputusan
diambil semata-mata dari keadaan yang teramati saat itu.

Tiga langkah berikut selalu ada pada agen mana pun, dari yang paling sederhana
sampai yang paling rumit.

\begin{langkah}
  \item \textbf{Persepsi.} Mengumpulkan keterangan tentang keadaan di
        sekitarnya. Pada bab ini keterangan itu hanya satu: di mana pemain
        berada sekarang.
  \item \textbf{Keputusan.} Menentukan tindakan yang dipilih. Pada bab ini
        hanya ada dua pilihan, yaitu mengejar atau berdiam.
  \item \textbf{Tindakan.} Melaksanakan keputusan itu dengan mengubah kecepatan,
        lalu posisi.
\end{langkah}

Yang menarik dari susunan itu adalah keluwesannya. Menambahkan perilaku baru
berarti menambahkan pilihan pada langkah kedua, sedangkan langkah pertama dan
ketiga tidak perlu berubah. Sifat itulah yang membuat susunan ini dapat
dikembangkan menjadi mesin keadaan pada Bab~IX, dan dilengkapi pencarian jalur
pada Bab~X.

\subsection{Perilaku yang Dibangun pada Bab Ini}

Perilaku yang dibangun sederhana dan sengaja demikian: NPC mengejar pemain
apabila pemain berada di dalam jangkauannya, dan berdiam apabila pemain berada
di luar jangkauan. Walaupun sederhana, perilaku ini sudah memuat seluruh
persoalan penting: bagaimana jarak dihitung, bagaimana arah ditentukan,
bagaimana kecepatan dibatasi, dan kapan berhenti.

Tiga hal berikut sengaja belum ada, dan masing-masing menuntut bab tersendiri.

\begin{itemize}
  \item \textbf{Keadaan yang lebih dari dua}, misalnya berpatroli dan
        menyerang. Dibahas pada Bab~IX.
  \item \textbf{Pengetahuan tentang rintangan}, sehingga NPC dapat mencari
        jalan memutar. Dibahas pada Bab~X.
  \item \textbf{Pertimbangan biaya perhitungan} ketika jumlah NPC bertambah.
        Dibahas pada Bab~XI.
\end{itemize}
